% p088_1; source PDF page 88; manually verified editable table.
\begin{tikzpicture}[x=1cm,y=-1cm,inner sep=0pt,outer sep=0pt]
\path[use as bounding box] (0,0) rectangle (16.000,5.130);
\draw[line width=.4pt] (0.000,0.000) rectangle (5.500,0.855);
\node[align=center,text width=5.240cm,font={\linespread{1}\fontsize{10}{13.50}\selectfont\bfseries}] at (2.750,0.427) {正当防卫的成立条件};
\draw[line width=.4pt] (5.500,0.000) rectangle (16.000,0.855);
\node[align=center,text width=10.240cm,font={\linespread{1}\fontsize{10}{13.50}\selectfont\bfseries}] at (10.750,0.427) {具体内容};
\draw[line width=.4pt] (0.000,0.855) rectangle (5.500,1.710);
\node[align=left,text width=5.240cm,font={\linespread{1}\fontsize{10}{13.50}\selectfont}] at (2.750,1.282) {起因条件};
\draw[line width=.4pt] (5.500,0.855) rectangle (16.000,1.710);
\node[align=left,text width=10.240cm,font={\linespread{1}\fontsize{10}{13.50}\selectfont}] at (10.750,1.282) {面临的侵害具有不法性、客观性和现实性（假想防卫）};
\draw[line width=.4pt] (0.000,1.710) rectangle (5.500,2.565);
\node[align=left,text width=5.240cm,font={\linespread{1}\fontsize{10}{13.50}\selectfont}] at (2.750,2.138) {时间条件};
\draw[line width=.4pt] (5.500,1.710) rectangle (16.000,2.565);
\node[align=left,text width=10.240cm,font={\linespread{1}\fontsize{10}{13.50}\selectfont}] at (10.750,2.138) {不法侵害正在进行（事后防卫）};
\draw[line width=.4pt] (0.000,2.565) rectangle (5.500,3.420);
\node[align=left,text width=5.240cm,font={\linespread{1}\fontsize{10}{13.50}\selectfont}] at (2.750,2.992) {意思条件};
\draw[line width=.4pt] (5.500,2.565) rectangle (16.000,3.420);
\node[align=left,text width=10.240cm,font={\linespread{1}\fontsize{10}{13.50}\selectfont}] at (10.750,2.992) {防卫者是否需要防卫意思（偶然防卫）};
\draw[line width=.4pt] (0.000,3.420) rectangle (5.500,4.275);
\node[align=left,text width=5.240cm,font={\linespread{1}\fontsize{10}{13.50}\selectfont}] at (2.750,3.848) {对象条件};
\draw[line width=.4pt] (5.500,3.420) rectangle (16.000,4.275);
\node[align=left,text width=10.240cm,font={\linespread{1}\fontsize{10}{13.50}\selectfont}] at (10.750,3.848) {防卫手段针对不法侵害人本人};
\draw[line width=.4pt] (0.000,4.275) rectangle (5.500,5.130);
\node[align=left,text width=5.240cm,font={\linespread{1}\fontsize{10}{13.50}\selectfont}] at (2.750,4.703) {限度条件};
\draw[line width=.4pt] (5.500,4.275) rectangle (16.000,5.130);
\node[align=left,text width=10.240cm,font={\linespread{1}\fontsize{10}{13.50}\selectfont}] at (10.750,4.703) {防卫手段具有必要性和相当性（防卫过当）};
\end{tikzpicture}
